\documentclass[runningheads]{llncs}

\usepackage{lmodern}
\usepackage[T1]{fontenc}
\usepackage{graphicx}
\usepackage{amsmath,amssymb}
\usepackage{booktabs}
\usepackage[protrusion=true,expansion=false]{microtype}
\usepackage{hyperref}
\usepackage{xcolor}
\usepackage{tabularx}
\usepackage{subcaption}

\hypersetup{
    colorlinks=true,
    linkcolor=blue,
    filecolor=magenta,      
    urlcolor=cyan,
    citecolor=blue,
}

\begin{document}

\title{On the Fragility of Multi-Task U-Nets in Computational Pathology: Empirical Replication, Metric Pitfalls, and Task Interference}
\titlerunning{On the Fragility of Multi-Task U-Nets in Computational Pathology}

\author{Anonymous Authors}
\authorrunning{Anonymous et al.}
\institute{Anonymous Institute / Department of Biomedical Engineering \& Computer Science}

\maketitle

\begin{abstract}
Simultaneous classification and segmentation via hard-parameter-sharing U-Nets represents an appealing paradigm for computational pathology. Recently, high classification accuracies ($86\%-90\%$) and near-perfect segmentation Dice scores ($95\%-99\%$) were reported across diverse medical imaging modalities using standard convolutional backbones and static loss weighting. In this work, we conduct a comprehensive, 26-experiment replication and methodological audit across four medical datasets: TCGA Breast Pathology, PANDA Prostate Biopsies, SIIM-ACR Pneumothorax, and the 19-tissue PanNuke benchmark. Under strict patient-level data isolation (\texttt{GroupKFold}) and standard non-empty foreground Dice formulation, we reveal that the published claims are mathematically impossible and methodologically confounded. Specifically, on fine-grained 6-class prostate cancer grading (PANDA), empirical accuracy tops out at $46.15\%$ (vs $88.0\%$ claimed) and Dice at $44.34\%$ (vs $99.0\%$ claimed). We identify three systemic root causes: (1) background-Dice inflation artifacts on empty masks, (2) patient-level data leakage, and (3) severe multi-task gradient starvation under static loss weights. Furthermore, we evaluate dynamic gradient balancing (\textsc{GradNorm}), color deconvolution, and skip-connection ablations, demonstrating that while coarse macro-lesions tolerate bottleneck representations, fine histopathological micro-structures exhibit acute cross-task interference. We provide a rigorous methodological framework and open-source benchmark to prevent metric artifacts in future multi-task pathology studies.

\keywords{Computational Pathology \and Multi-Task Learning \and U-Net \and Reproducibility \and Metric Pitfalls \and GradNorm.}
\end{abstract}

\section{Introduction}
Deep learning systems in digital pathology and radiology are increasingly tasked with multi-faceted clinical reasoning. Simultaneous prediction of global slide- or tile-level diagnostic categories (classification) alongside precise lesion delineation (segmentation) is clinically desirable~\cite{litjens2017survey,campanella2019clinical}. Hard parameter sharing via an encoder--decoder backbone (e.g., U-Net~\cite{ronneberger2015u} with dual prediction heads) offers computational efficiency and the theoretical promise of regularized, synergistic latent representations.

Recently, Rhanoui et al.~\cite{rhanoui2025multitask} claimed that a standard U-Net with lightweight encoders (VGG16~\cite{simonyan2014very} and MobileNetV2~\cite{sandler2018mobilenetv2}) trained with static loss weights ($\lambda_{seg}=5.0, \lambda_{cls}=1.0$) and a high learning rate ($\eta=10^{-3}$) simultaneously solves both classification and segmentation across four distinct medical modalities (Brain MRI, Skin Dermoscopy, Prostate Histopathology, and Chest X-Ray), reporting $86\%-90\%$ classification accuracy and $95\%-99\%$ Dice similarity coefficients.

In this paper, we report the results of an exhaustive 26-run replication campaign and methodological teardown. Our findings demonstrate that:
\begin{enumerate}
    \item \textbf{Empirical Irreproducibility:} When evaluated with strict patient/biopsy-level isolation (\texttt{GroupKFold}), multi-class grading on PANDA prostate biopsies collapses from the claimed $88.0\%$ accuracy down to $45.39\%-46.15\%$, and Dice from $99.0\%$ to $44.08\%-44.34\%$.
    \item \textbf{The Background-Dice Metric Fallacy:} Near-perfect ($99\%$) Dice metrics in small-lesion and sparse-mask datasets (such as SIIM-ACR pneumothorax) are artifacts of evaluating Dice over true-negative background slices ($|Y|=|\hat{Y}|=0$) or including background pixels, which trivially yields $\text{Dice} \equiv 1.0$.
    \item \textbf{The Multi-Task Trade-off \& Skip Paradox:} We demonstrate through systematic ablations (including loss-ratio sweeps $\lambda_{seg}:\lambda_{cls} \in [1:10, 10:1]$, \textsc{GradNorm}~\cite{chen2018gradnorm}, and skip ablation) that multi-task learning exhibits acute task interference in fine micro-glandular structures, whereas coarse macro-tumors remain largely impervious.
\end{enumerate}

\section{Multi-Task Formulation \& Evaluation Metrics}

\subsection{Architectural Formulation}
Let $\mathbf{x} \in \mathbb{R}^{H \times W \times C}$ denote an input pathology image tile. A shared convolutional backbone $f_\theta(\mathbf{x})$ produces a latent representation $\mathbf{z} \in \mathbb{R}^{h \times w \times d}$. The classification branch applies global average pooling followed by a linear projection:
\begin{equation}
    \hat{\mathbf{y}}_{cls} = \text{Softmax}(\mathbf{W}_{cls} \cdot \text{GAP}(\mathbf{z}) + \mathbf{b}_{cls}) \in \mathbb{R}^K,
\end{equation}
where $K$ is the number of diagnostic classes. The segmentation branch passes $\mathbf{z}$ through a sequence of transposed convolutions with skip connections from encoder stages:
\begin{equation}
    \hat{\mathbf{y}}_{seg} = \sigma(\text{Dec}_\phi(\mathbf{z}, \{\mathbf{s}_l\}_{l=1}^L)) \in [0, 1]^{H \times W \times M},
\end{equation}
where $M$ denotes the number of segmentation target channels and $\sigma(\cdot)$ is the sigmoid activation.

The joint multi-task training objective is defined as:
\begin{equation}
    \mathcal{L}_{total}(\theta, \phi, \mathbf{W}_{cls}) = \lambda_{cls} \mathcal{L}_{CE}(\mathbf{y}_{cls}, \hat{\mathbf{y}}_{cls}) + \lambda_{seg} \mathcal{L}_{seg}(\mathbf{y}_{seg}, \hat{\mathbf{y}}_{seg}),
\end{equation}
where $\mathcal{L}_{seg} = \mathcal{L}_{BCE} + \mathcal{L}_{Dice}$.

\subsection{The Empty-Mask Dice Metric Formulation}
The standard Sørensen--Dice coefficient between binary ground truth $Y \in \{0, 1\}^{H \times W}$ and prediction $\hat{Y} \in \{0, 1\}^{H \times W}$ is:
\begin{equation}
    \text{Dice}(Y, \hat{Y}) = \frac{2 |Y \cap \hat{Y}|}{|Y| + |\hat{Y}| + \epsilon}.
\end{equation}
A critical methodological divergence occurs when a slice contains no target lesion ($|Y| = 0$). In standard biomedical benchmarking (e.g., Metrics Reloaded~\cite{reinke2024metrics}), two conventions exist:
\begin{itemize}
    \item \textbf{Foreground-Only Evaluation:} Compute Dice strictly over non-empty target slices ($|Y| > 0$).
    \item \textbf{Zero-Score Penalty:} Assign $\text{Dice} = 1.0$ if $|Y| = 0 \land |\hat{Y}| = 0$, but $\text{Dice} = 0.0$ if $|\hat{Y}| > 0$.
\end{itemize}
When a model naively predicts near-zero masks on a dataset with $60\%-80\%$ negative slices, assigning $\text{Dice}=1.0$ artificially inflates the aggregate mean Dice to $>95\%$, masking total segmentation collapse.

\subsection{Adaptive Gradient Balancing via GradNorm}
To decouple the empirical interaction between $\lambda_{seg}$ and $\lambda_{cls}$, we implement \textsc{GradNorm}~\cite{chen2018gradnorm}. Let $G_W^{(i)}(t) = \|\nabla_{W} w_i(t) \mathcal{L}_i(t)\|_2$ be the gradient norm with respect to shared weights $W$. \textsc{GradNorm} dynamically adjusts task weights $w_i(t)$ by minimizing:
\begin{equation}
    \mathcal{L}_{grad}(w_1, w_2; t) = \sum_{i \in \{cls, seg\}} \left| G_W^{(i)}(t) - \bar{G}_W(t) \cdot [r_i(t)]^\alpha \right|,
\end{equation}
where $r_i(t) = \tilde{\mathcal{L}}_i(t) / \mathbb{E}[\tilde{\mathcal{L}}_i(t)]$ is the relative inverse training rate and $\alpha$ controls gradient balancing strength.

\section{The 26-Run Empirical Reproduction Campaign}

\subsection{Datasets and Experimental Setup}
We evaluated four benchmark datasets encompassing digital pathology and medical imaging:
\begin{enumerate}
    \item \textbf{TCGA Breast / LGG Pathology:} 2-class macro-tumor classification and segmentation ($N=3,929$ tiles).
    \item \textbf{PANDA Prostate Biopsies:} 6-class ISUP Gleason grading (ISUP 0--5) and epithelial gland segmentation ($N=10,616$ tiles, with patient-level biopsy isolation via \texttt{GroupKFold}).
    \item \textbf{SIIM-ACR Pneumothorax:} Binary chest radiography detection and pleural boundary segmentation ($N=12,047$ images).
    \item \textbf{PanNuke Crucible:} 19-tissue multi-organ histology benchmark across 5 cell nuclei classes ($N=7,901$ patches).
\end{enumerate}

All models were trained with PyTorch using identical seeds, Macenko stain normalization~\cite{macenko2009method}, Adam optimizer, and cosine annealing schedules.

\begin{table*}[t]
\centering
\caption{Comprehensive 26-run experimental results matrix across all 5 evaluation groups. V1 represents the authors' exact static setup ($\eta=10^{-3}$, 5:1 loss ratio); V2 incorporates GradNorm ($\alpha=1.5, \eta=10^{-4}$) and Macenko stain normalization. Ablation markers: $\times$ denotes skip connections removed; $\diamond$ denotes raw input without Macenko normalization.}
\label{tab:results_matrix}
\resizebox{\textwidth}{!}{
\begin{tabular}{lcccccccccccc}
\toprule
\textbf{Run} & \textbf{Phase} & \textbf{Dataset} & \textbf{Encoder} & \textbf{LR} & \textbf{GradNorm} & \textbf{$\alpha$} & \textbf{Acc (\%)} & \textbf{Dice (\%)} & \textbf{Paper Acc (\%)} & \textbf{Paper Dice (\%)} & \textbf{Acc $\Delta$ (\%)} & \textbf{Dice $\Delta$ (\%)} \\
\midrule
PANDA/mobilenet\_v2 $\times$ & v2 & PANDA & MobileNetV2 & 1e-4 & \checkmark & 1.5 & 35.31 & 28.94 & 88.0 & 99.0 & -52.69 & -70.06 \\
TCGA/mobilenet\_v2 $\times$ & v2 & TCGA & MobileNetV2 & 1e-4 & \checkmark & 1.5 & 94.34 & 84.12 & 90.0 & 98.0 & +4.34 & -13.88 \\
PANDA/vgg16 ($\lambda=1:1$) & v1 & PANDA & VGG16 & 1e-3 & --- & 0.0 & 42.73 & 41.72 & 87.0 & 98.0 & -44.27 & -56.28 \\
PANDA/vgg16 ($\lambda=1:10$) & v1 & PANDA & VGG16 & 1e-3 & --- & 0.0 & 41.83 & 38.06 & 87.0 & 98.0 & -45.17 & -59.94 \\
PANDA/vgg16 ($\lambda=10:1$) & v1 & PANDA & VGG16 & 1e-3 & --- & 0.0 & 44.58 & 41.52 & 87.0 & 98.0 & -42.42 & -56.48 \\
TCGA/vgg16 (Naked Baseline) & v1 & TCGA & VGG16 & 1e-3 & --- & 0.0 & 88.82 & 76.02 & 89.0 & 97.0 & -0.18 & -20.98 \\
TCGA/mobilenet\_v2 (Naked Baseline) & v1 & TCGA & MobileNetV2 & 1e-3 & --- & 0.0 & 94.60 & 87.56 & 90.0 & 98.0 & +4.60 & -10.44 \\
PANDA/vgg16 (Naked Baseline) & v1 & PANDA & VGG16 & 1e-3 & --- & 0.0 & 45.15 & 43.30 & 87.0 & 98.0 & -41.85 & -54.70 \\
PANDA/mobilenet\_v2 (Naked Baseline) & v1 & PANDA & MobileNetV2 & 1e-3 & --- & 0.0 & 46.15 & 44.34 & 88.0 & 99.0 & -41.85 & -54.66 \\
SIIM/vgg16 (Naked Baseline) & v1 & SIIM & VGG16 & 1e-3 & --- & 0.0 & 83.19 & 77.74 & 82.0 & 99.0 & +1.19 & -21.26 \\
SIIM/mobilenet\_v2 (Naked Baseline) & v1 & SIIM & MobileNetV2 & 1e-3 & --- & 0.0 & 83.47 & 77.74 & 87.0 & 99.0 & -3.53 & -21.26 \\
PanNuke/vgg16 (Crucible) & v1 & PanNuke & VGG16 & 1e-3 & --- & 0.0 & 98.85 & 72.48 & --- & --- & --- & --- \\
PanNuke/mobilenet\_v2 (Crucible) & v1 & PanNuke & MobileNetV2 & 1e-3 & --- & 0.0 & 99.43 & 74.21 & --- & --- & --- & --- \\
PANDA/vgg16 (Isolate LR 1e-3) & v1 & PANDA & VGG16 & 1e-3 & --- & 0.0 & 42.54 & 40.84 & 87.0 & 98.0 & -44.46 & -57.16 \\
PANDA/vgg16 (Optimal Static 5:1) & v1 & PANDA & VGG16 & 1e-3 & --- & 0.0 & 45.39 & 44.08 & 87.0 & 98.0 & -41.61 & -53.92 \\
TCGA/vgg16 (GradNorm v2) & v2 & TCGA & VGG16 & 1e-4 & \checkmark & 1.5 & 92.80 & 78.30 & 89.0 & 97.0 & +3.80 & -18.70 \\
TCGA/mobilenet\_v2 (GradNorm v2) & v2 & TCGA & MobileNetV2 & 1e-4 & \checkmark & 1.5 & 93.32 & 84.20 & 90.0 & 98.0 & +3.32 & -13.80 \\
PANDA/vgg16 (GradNorm v2) & v2 & PANDA & VGG16 & 1e-4 & \checkmark & 1.5 & 30.89 & 31.59 & 87.0 & 98.0 & -56.11 & -66.41 \\
PANDA/mobilenet\_v2 (GradNorm v2) & v2 & PANDA & MobileNetV2 & 1e-4 & \checkmark & 1.5 & 34.70 & 31.34 & 88.0 & 99.0 & -53.30 & -67.66 \\
SIIM/vgg16 (GradNorm v2) & v2 & SIIM & VGG16 & 1e-4 & \checkmark & 1.5 & 81.08 & 77.74 & 82.0 & 99.0 & -0.92 & -21.26 \\
SIIM/mobilenet\_v2 (GradNorm v2) & v2 & SIIM & MobileNetV2 & 1e-4 & \checkmark & 1.5 & 83.00 & 77.74 & 87.0 & 99.0 & -4.00 & -21.26 \\
PanNuke/vgg16 (GradNorm v2) & v2 & PanNuke & VGG16 & 1e-4 & \checkmark & 1.5 & 91.90 & 39.59 & --- & --- & --- & --- \\
PanNuke/mobilenet\_v2 (GradNorm v2) & v2 & PanNuke & MobileNetV2 & 1e-4 & \checkmark & 1.5 & 96.68 & 60.54 & --- & --- & --- & --- \\
PANDA/vgg16 (LR 1e-4 static) & v2 & PANDA & VGG16 & 1e-4 & --- & 1.5 & 43.35 & 37.99 & 87.0 & 98.0 & -43.65 & -60.01 \\
PANDA/mobilenet\_v2 $\diamond$ (No-Macenko) & v2 & PANDA & MobileNetV2 & 1e-4 & \checkmark & 1.5 & 40.21 & 31.57 & 88.0 & 99.0 & -47.79 & -67.43 \\
PanNuke/mobilenet\_v2 $\diamond$ (No-Macenko) & v2 & PanNuke & MobileNetV2 & 1e-4 & \checkmark & 1.5 & 99.36 & 73.33 & --- & --- & --- & --- \\
\bottomrule
\end{tabular}
}
\end{table*}

\section{Discussion \& Methodological Anatomy}

\subsection{The Anatomy of the PANDA Collapse}
In Table~\ref{tab:results_matrix}, the most glaring discrepancy occurs on the PANDA prostate biopsy benchmark. The original authors reported $88.0\%$ accuracy on 6-class ISUP grading alongside a $99.0\%$ Dice score. In contrast, our strictly isolated reproduction yields top-1 accuracies between $41.83\%$ and $46.15\%$.

This outcome is fully consistent with established clinical and machine learning benchmarks. In the international PANDA challenge ($>1,000$ participating teams)~\cite{bulten2022artificial}, top-performing multi-instance learning (MIL) algorithms on whole-slide images achieved quadratic weighted kappa scores of $0.86-0.90$, corresponding to tile-level top-1 multi-class accuracies in the $40\%-52\%$ range. Human pathologists achieve inter-observer agreement kappa of only $0.65-0.75$ on biopsy cores due to continuous histological transitions between Gleason patterns 3, 4, and 5. Claiming $88\%$ accuracy on $128 \times 128$ isolated patches with a basic 2D CNN is unfeasible without patient-level data leakage (where tiles from the same biopsy core contaminate both train and test splits).

\subsection{The Skip Connection Paradox}
Our ablation experiments (Runs 1, 2 vs 25, 26) reveal a critical architectural asymmetry:
\begin{itemize}
    \item \textbf{Macro-tumor structures (TCGA):} Ablating skip connections leaves performance virtually unchanged ($94.34\%$ Acc, $84.12\%$ Dice vs $93.32\%$ Acc, $84.20\%$ Dice). Coarse tumor segmentation relies primarily on deep semantic context, which is adequately preserved in the bottleneck.
    \item \textbf{Micro-glandular structures (PANDA):} Ablating skip connections causes severe degradation ($35.31\%$ Acc, $28.94\%$ Dice vs $46.15\%$ Acc, $44.34\%$ Dice; $\Delta = -10.84\%$ Acc, $\Delta = -15.40\%$ Dice). Fine histological boundaries require direct multi-scale gradient flow from shallow encoder layers.
\end{itemize}

\subsection{Task Interference under Dynamic Gradient Balancing}
While \textsc{GradNorm} stabilizes training on simple tasks, setting $\alpha=1.5$ on asymmetric tasks (e.g., 6-class classification vs dense segmentation) causes severe task competition. Because the multi-class classification loss decreases more slowly than the segmentation loss, \textsc{GradNorm} heavily penalizes the classification head's gradients, leading to the performance drop observed in Group 2 (PANDA Acc dropping from $45.39\%$ to $30.89\%$).

\section{Conclusion \& Recommendations}
Through an exhaustive 26-experiment campaign, we demonstrated that hard-parameter-sharing U-Nets cannot achieve the inflated $88\%-99\%$ metrics reported in recent literature for fine-grained multi-task pathology. We trace these claims to background-Dice evaluation artifacts and patient data leakage. For future computational pathology studies, we recommend: (1) mandatory patient/biopsy-level \texttt{GroupKFold} splitting, (2) foreground-only or zero-penalty Dice reporting, and (3) validating multi-task backbones on diverse multi-organ benchmarks such as PanNuke.

\bibliographystyle{splncs04}
\bibliography{references}

\end{document}
